% TFM-ML-cap02-marco-contextual.tex
% Capítulo 02 — Marco Contextual
% Autor: Mario Lourido Regueira
% Scope: v5 — Corredor CHUF–Ferrol, horizonte variable
% Última revisión: 2026-04-30
% Estado: REDACTADO — §2.1, §2.2 y §2.3 alineados con v5
%         §2.4–2.5 (Caso Ferrol y Stakeholders) en espera de verificación


% ============================================================
% 2.1 EL VACÍO DE TRANSPORTE NOCTURNO EN EUROPA
% ============================================================

\subsection{El vacío de transporte nocturno en Europa}
\label{sec:vacio-nocturno}

Las ciudades europeas funcionan las veinticuatro horas del día. Los
hospitales no cierran, las cadenas de distribución no se detienen, los
hoteles no apagan sus luces. Sin embargo, los sistemas de transporte
público que conectan a sus trabajadores con sus lugares de trabajo sí lo
hacen, de forma sistemática, durante las horas de menor rentabilidad
comercial: la madrugada.

Esta discontinuidad responde a una lógica
económica que prima la eficiencia de red sobre la cobertura de necesidades:
a partir de cierta hora, la demanda cae por debajo del umbral de rentabilidad
del servicio convencional y el operador retira los vehículos. El resultado
es un vacío de entre cuatro y seis horas en el que los trabajadores de turno
nocturno deben resolver su movilidad por medios propios (vehículo
privado, taxi o, en los casos más vulnerables, a pie).

\begin{claimS}
Ese vacío se concentra típicamente entre la 01:00 y las 05:00~h
\parencite{Plyushteva2020}. A principios de 2019, 26 de las 28 capitales de
la Unión Europea contaban con alguna forma de transporte público nocturno,
aunque con frecuencias y coberturas muy reducidas respecto al servicio diurno
\parencite{Plyushteva2020}. Las ciudades no capitales (y especialmente
aquellas de menos de 150.000~habitantes) quedan en su mayor parte fuera
de este dato: sus redes nocturnas son, cuando existen, meramente testimoniales.
\end{claimS}

El vacío nocturno no es un problema homogéneo. Afecta de forma diferenciada
según tres variables: el perfil del trabajador (turno fijo o rotativo,
distancia al centro de trabajo, nivel de renta), el contexto geográfico
(ciudad grande con red nocturna deficiente frente a ciudad media sin red
alguna) y el género (la percepción de inseguridad nocturna actúa como
barrera adicional para las mujeres, que constituyen la mayoría del empleo
sanitario nocturno). El diseño del servicio propuesto en este trabajo trata
las tres variables de forma explícita.


% ============================================================
% 2.2 PANORAMA DE SOLUCIONES DE MOVILIDAD NOCTURNA
% ============================================================

\subsection{Panorama de soluciones de movilidad nocturna: modelos operativos existentes}
\label{sec:panorama-soluciones}

Antes de formular ninguna propuesta, es necesario describir
qué soluciones existen para cubrir vacíos de transporte
nocturno en corredores de demanda concentrada. Esta sección
presenta los modelos operativos documentados en la
literatura y en experiencias europeas. Ninguno se selecciona
aquí: la evaluación frente a los criterios de diseño
se realiza en el capítulo~\ref{cap:propuesta}.

\subsubsection{Transporte público convencional ampliado}
\label{sec:sol-tp-convencional}

La extensión de horarios de líneas de autobús o tren
existentes es la solución más habitual para cubrir
parcialmente la franja nocturna. Su viabilidad depende
de la existencia de demanda suficiente para justificar
el coste operativo. En ciudades de menos de 150.000
habitantes, la demanda nocturna dispersa raramente
alcanza el umbral de rentabilidad de una línea regular
\parencite{Estrada2021}.

\subsubsection{Transporte a demanda (DRT)}
\label{sec:sol-drt}

El \emph{Demand Responsive Transport} (DRT) adapta rutas
y horarios a las solicitudes de los usuarios en tiempo
real o con reserva anticipada. Reduce el coste por
kilómetro en contextos de baja densidad de demanda y
permite cubrir la franja nocturna sin la rigidez de
una línea fija. Varias ciudades europeas de tamaño medio
han implementado DRT nocturno con resultados variables
según el modelo de financiación y la tasa de adopción
\parencite{Busitalia2019QuiBus, K2Centrum2025DRT}.

\subsubsection{Shuttle con conductor en corredor fijo}
\label{sec:sol-shuttle-conductor}

Un shuttle de capacidad reducida (8--20 plazas) operando
en un corredor fijo y con horarios sincronizados con
un generador de demanda concentrada (hospital, polígono
industrial, campus logístico) es una solución de baja
complejidad tecnológica y alta adaptabilidad. No requiere
infraestructura específica ni marco regulatorio especial.
Su coste operativo está determinado principalmente por
el salario del conductor y la amortización del vehículo
\parencite{Stagecoach2020Hull}.

\subsubsection{Vehículo autónomo en corredor fijo}
\label{sec:sol-av}

Los vehículos autónomos de nivel SAE~L4 operando en
corredor definido y geofenced eliminan el coste del
conductor a largo plazo y permiten operación continua
sin límites de jornada laboral. Los pilotos europeos
más relevantes en entorno urbano (SHOW H2020, AVENUE H2020)
han demostrado viabilidad técnica en condiciones
controladas, con tasas de aceptación de usuario variables
según el contexto y la franja horaria
\parencite{SHOW2024Legacy}.
Las limitaciones actuales incluyen la madurez regulatoria,
el rendimiento en condiciones meteorológicas adversas
y el coste de la infraestructura de mapeo previo.

\subsubsection{Ride-hailing y acuerdos institucionales}
\label{sec:sol-ridehailing}

Las plataformas de \emph{ride-hailing} (VTC, taxi) y los
acuerdos de transporte institucional entre el empleador
y un operador externo representan soluciones de implantación
rápida sin necesidad de flota propia. Su coste por viaje
es generalmente superior al de un servicio de línea, pero
puede ser asumido total o parcialmente por la institución
empleadora como beneficio social \parencite{Stagecoach2020Hull, K2Centrum2025DRT}.

\subsubsection{Síntesis comparativa}
\label{sec:sol-sintesis}

Cada modelo presenta un perfil distinto de coste, madurez,
complejidad de implantación y cobertura de los criterios
de diseño que se establecerán en el capítulo~\ref{cap:criterios}.
La evaluación sistemática de estos candidatos frente a
dichos criterios se realiza en §\ref{sec:evaluacion-candidatos}.


% ============================================================
% SECCIONES PENDIENTES (cap02)
% ============================================================

% ============================================================
%  §2.3 ANÁLISIS PESTEL
%  Archivo: TFM-ML-cap02-marco-contextual.tex
%  Estado: primera redacción — pendiente revisión Mario
%  Claves .bib verificadas en sesión 2026-03-13
% ============================================================

\subsection{Análisis PESTEL: macroentorno de la movilidad nocturna en ciudades medias europeas}
\label{sec:pestel}

El análisis PESTEL permite enmarcar las condiciones de viabilidad del
servicio propuesto más allá del corredor concreto de Ferrol. Se analiza
el macroentorno de la movilidad nocturna en Europa a lo largo de seis dimensiones, con énfasis en los factores que determinan si las soluciones de transporte pueden operar de forma sostenida en ciudades de tamaño medio con vacíos de transporte nocturno documentados.

% ------------------------------------------------------------------
\subsubsection{Dimensión política}
\label{sec:pestel-politica}

El marco regulatorio europeo para la conducción automatizada ha alcanzado
un grado de madurez institucional que convierte el horizonte 2035 en
plausible. El Reglamento de Ejecución~(UE) 2022/1426 de la Comisión
estableció por primera vez los procedimientos uniformes y las
especificaciones técnicas para la homologación de tipo del sistema de
conducción automatizada~(ADS) de vehículos totalmente automatizados,
fijando las condiciones bajo las cuales los fabricantes pueden obtener
homologación en la Unión~\parencite{RegUE20221426}. Este marco fue
actualizado en marzo de 2026 mediante el Reglamento de Ejecución~(UE)
2026/481, cuyo preámbulo reconoce explícitamente que los avances
tecnológicos y regulatorios acumulados desde 2022 hacen posible el
despliegue a gran escala de vehículos totalmente
automatizados~\parencite{RegUE20262026481}. La actualización es
significativa: supone que la UE deja de tratar el vehículo autónomo como
un objeto de regulación experimental para integrarlo en el sistema de
homologación ordinario.

En el plano nacional, la Dirección General de Tráfico articula el
programa marco~ES-AV como instrumento de autorización de pruebas de
vehículos automatizados en España, proporcionando el vínculo operativo
entre el reglamento europeo y los actores del ecosistema en
territorio~español~\parencite{DGTESAV2026}. En paralelo, la iniciativa
URBACT \textit{Cities After Dark} visibiliza la movilidad nocturna como
objeto de política pública municipal, abriendo la vía para que las
ciudades de tamaño medio incorporen la franja nocturna en sus Planes de
Movilidad Urbana Sostenible~(SUMP)~\parencite{URBACT2025Gender}. Ambos
instrumentos (europeo y local) convergen en el tipo de servicio que este
proyecto propone: operación en corredor fijo, entorno regulado y demanda
concentrada y predecible. Este marco regulatorio es relevante si la solución
seleccionada en la fase de diseño implica niveles de automatización que lo requieran.

% ------------------------------------------------------------------
\subsubsection{Dimensión económica}
\label{sec:pestel-economica}

La economía del shuttle autónomo nocturno presenta una estructura de
costes con barreras iniciales elevadas y ventajas comparativas
crecientes respecto a las alternativas actuales. El cartografiado de alta
definición~(HD~Maps) necesario para la operación de un ADS en entorno
urbano supone entre 5.000 y 25.000~euros por kilómetro de carretera
mapeada, dependiendo de la densidad del entorno y la resolución
requerida~\parencite{Moreau2023}. Para un corredor de entre
cinco y ocho kilómetros como el del corredor~CHUF, este coste es
perfectamente absorbible dentro del presupuesto de un proyecto piloto,
precisamente porque la extensión reducida del corredor es una ventaja
estructural del modelo que este TFM propone.

La demanda latente existe y tiene un coste de oportunidad medible. El
estudio de~\textcite{PlyushetvaBoussauw2020} sobre trabajadores de turno
nocturno en Sofía documenta que un empleador del sector llegaba a asumir
más de 2.000~levs mensuales (+-1020€) en taxis nocturnos para garantizar la
asistencia del personal, cifra que evidencia que la disposición a pagar
(tanto por parte de empleadores como de instituciones) es una práctica documentada. 
La disposición a pagar de los usuarios
finales es igualmente relevante: estudios de aceptación del vehículo
autónomo compartido apuntan a que los usuarios de transporte público en
entorno urbano asignan valor monetario significativo a las prestaciones
de~L4, especialmente cuando la percepción de seguridad nocturna es baja
\parencite{Gavanas2024}. La implicación para el diseño del servicio es que
el modelo no necesita competir en precio con el transporte público diurno,
sino con el taxi nocturno (comparación que mejora sustancialmente la
ecuación de valor).

% ------------------------------------------------------------------
\subsubsection{Dimensión social}
\label{sec:pestel-social}

La dimensión social del servicio está determinada por dos realidades que
se superponen: la composición de género del colectivo objetivo y la
percepción diferencial de seguridad nocturna. El empleo en el sector
sanitario y de trabajo social en la Unión Europea es mayoritariamente
femenino: según datos de Eurostat correspondientes al tercer trimestre
de~2021, el~78\% de los trabajadores de este sector son
mujeres~\parencite{Eurostat2022health}. Esta cifra es un dato estructural: 
cualquier servicio diseñado para el personal
sanitario nocturno tiene a las mujeres como usuaria principal, lo que
convierte la perspectiva de género en un requisito de diseño.

La percepción de seguridad nocturna a su vez muestra una brecha de género
cuantificada. \textcite{PlyushetvaBoussauw2020} midieron en Sofía que
los hombres valoraban su seguridad nocturna en 5,82 sobre~10 frente a
4,42 de las mujeres, diferencia que se traduce directamente en menor
disposición a utilizar el transporte público durante la franja nocturna y
en mayor dependencia de alternativas privadas de mayor coste. El diseño
del espacio interior del shuttle, su iluminación y sus mecanismos de
comunicación~eHMI deben responder a esta percepción diferencial, poniendola en valor.
En cuanto a la aceptación social de la tecnología autónoma en
sí, los pilotos del proyecto~SHOW muestran resultados favorables: los
usuarios de transporte público valoraron su experiencia en el shuttle
autónomo por encima de~7 en una escala~Likert de~9
puntos~\parencite{SHOW2024Legacy}, lo que indica que la barrera de
aceptación tecnológica es superable cuando el servicio se articula en
torno a necesidades reales de movilidad.

% ------------------------------------------------------------------
\subsubsection{Dimensión tecnológica}
\label{sec:pestel-tecnologica}

La dimensión tecnológica del macroentorno incluye tanto las soluciones de transporte convencional (digitalización de la gestión de flota, plataformas DRT) como las tecnologías emergentes de automatización del vehículo. Estas últimas se describen como factor del entorno, no como solución predeterminada.

La tecnología de conducción automatizada a nivel SAE~L4 existe en forma
de prototipo operativo, pero su despliegue comercial estable plantea
retos técnicos no resueltos que los propios pilotos europeos han
documentado. La operación nocturna introduce condiciones adversas
específicas: los pilotos del proyecto~SHOW identificaron en el municipio
de~Les~Mureaux~(Francia) que la iluminación deficiente genera puntos
ciegos para los sistemas de percepción del vehículo y obliga a solicitar
asistencia operativa externa con una frecuencia significativamente mayor
que en condiciones
diurnas~\parencite{SHOW2024Legacy}. En la misma línea, el
mantenimiento inadecuado de las marcas viales (una realidad habitual en
entornos periféricos como el corredor~CHUF) dificulta la detección por
cámara de los bordes de carril~\parencite{SHOWd812020}. Los mapas de alta
definición~(HD~Maps) utilizados por los sistemas ADS requieren una
precisión inferior a~10~cm para garantizar la localización
~\parencite{SHOWd812020}, lo que implica que cualquier cambio
en la vía (obra, repintado, nuevo obstáculo) debe reflejarse en tiempo
real en los mapas del sistema.

La madurez tecnológica general para el despliegue comercial a gran escala
no ha llegado aún en la actualidad inmediata (2026). \textcite{KonstantasAVENUE2023lessons} concluye en los
materiales de lecciones aprendidas del proyecto~AVENUE~H2020 que no se
esperan despliegues a gran escala de transporte autónomo durante al menos
cinco años desde el cierre del proyecto en~2022, lo que sitúa ese umbral
en torno a~2027–2028. Este horizonte no invalida el modelo propuesto para
2035: el corredor de Ferrol se plantea
como operación en entorno acotado y controlado, precisamente el formato
que los~\textit{stakeholders} del sector contemplan como viable en el
próximo decenio~\parencite{DibaGore2025}. El accidente del~Digibus en
Austria en agosto de~2023, que derivó en la revocación temporal del
permiso de pruebas, ilustra que los sistemas de supervisión remota y los
protocolos de intervención de emergencia son esenciales y
constitutivos dentro del
servicio~\parencite{SHOW2024Legacy}.

% ------------------------------------------------------------------
\subsubsection{Dimensión ecológica}
\label{sec:pestel-ecologica}

La operación de un shuttle eléctrico compartido en franja nocturna encaja
directamente en el marco del Pacto Verde Europeo, cuyo objetivo de
neutralidad climática para~2050 exige una transformación profunda del
modelo de movilidad urbana~\parencite{UITP2024}. El punto de partida no
es neutro: el transporte colectivo ya es intrínsecamente más eficiente
que el vehículo privado (el autobús duplica la eficiencia energética por
viajero-kilómetro del automóvil, y el ferrocarril la cuadriplica), ventaja que se acumula 
cuando la comparación se hace respecto al taxi
nocturno individual, alternativa modal predominante en ausencia de
servicio público nocturno~\parencite{UITP2024}. Los pilotos del proyecto
SHOW operan con flotas~100\% eléctricas~\parencite{SHOWd812020}, lo que
demuestra que la integración de la propulsión eléctrica en el shuttle
autónomo urbano es una opción técnica ya validada
en entorno de piloto. Para un corredor corto y de trayecto predecible
como el de~CHUF, la electrificación plena de la flota no enfrenta
limitaciones de autonomía de batería y permite planificar la recarga
sincronizada con los períodos de baja demanda entre cambios de turno.

% ------------------------------------------------------------------
\subsubsection{Dimensión legal}
\label{sec:pestel-legal}

El principal vacío regulatorio que afecta a la viabilidad del servicio
propuesto no es la homologación del vehículo (resuelta en términos
generales por el Reglamento~2022/1426 y su actualización de~2026) sino
la asignación de responsabilidad civil en caso de accidente. La literatura
especializada coincide en señalar que no existe un marco legal
comprehensivo en la~UE que resuelva con claridad a quién corresponde la
responsabilidad en un siniestro protagonizado por un vehículo autónomo:
¿al fabricante del sistema~ADS, al operador del servicio o al propietario
del vehículo?~\parencite{HafizAlKhafagy2025}. La Directiva~(UE)~2024/2853
sobre responsabilidad por productos defectuosos moderniza el régimen
existente extendiendo el concepto de «producto defectuoso» al
software~y~a~la~inteligencia~artificial~\parencite{DirectivaUE20242853},
lo que abre la vía para imputar responsabilidad al fabricante del~ADS en
caso de fallo del sistema. Sin embargo, la retirada de la propuesta de
Directiva de Responsabilidad en materia de~IA en~2025 por falta de
consenso político confirma que el marco de responsabilidad civil para
escenarios complejos de conducción automatizada sigue sin estar
armonizado en el conjunto de la~Unión.

En paralelo, el marco laboral europeo ofrece una palanca regulatoria en
favor del servicio propuesto: la Directiva~2003/88/CE establece
protecciones específicas para los trabajadores nocturnos, reconociendo el
carácter estructuralmente más oneroso de la jornada
nocturna~\parencite{Directiva200388CE}. Esta protección normativa
refuerza el argumento de que las instituciones públicas (SERGAS, Concello
de~Ferrol, Xunta de~Galicia) tienen fundamento legal para priorizar la
movilidad del personal de turno nocturno dentro de sus políticas de
bienestar laboral. La figura de las Obligaciones de Servicio
Público~(OSP) en el transporte colectivo permite estructurar el servicio
con garantías de viabilidad económica aun cuando la demanda nocturna no
alcance rentabilidad comercial autónoma~\parencite{UITP2024}. En España,
el programa~ES-AV de la~DGT establece el marco de autorización que
permite operar pruebas y servicios de vehículos automatizados en vía
pública bajo condiciones reguladas~\parencite{DGTESAV2026}.

\subsection{Caso de estudio: Corredor CHUF--Ferrol}
\label{sec:caso-ferrol}

\begin{pendiente}
%% INSERTAR: descripcion-corredor-ferrol

Ferrol es un municipio gallego situado en el noroeste de España de 64.218~habitantes según el Padrón Municipal a 1~de enero de 2024 \parencite{INE2024Ferrol}.

El \ac{CHUF}, que integra los hospitales Arquitecto Marcide, Naval y Profesor Novoa Santos, se ubica en la Avenida de la Residencia s/n, Ferrol \parencite{Sanidad2024CHUF}, fuera de la cuadrícula urbana central de la ciudad, a una distancia por carretera de entre 4{,}2 y
  5{,}2~km del centro urbano (Plaza de Armas),
  con un tiempo de trayecto de aproximadamente
  11~minutos en condiciones normales de tráfico
  \parencite{GoogleMaps2026Ferrol}.

Las consultas realizadas en los portales oficiales del Concello de Ferrol y del operador de transporte urbano no arrojan ninguna línea nocturna regular en la franja 22:00--06:00 en el corredor hospital--centro.
\end{pendiente}

\subsection{Mapa de stakeholders}
\label{sec:stakeholders}

\begin{pendiente}
Actores del ecosistema: operadores de transporte (Concello Ferrol, Xunta,
operadores privados), \textsc{sergas}, proveedores del servicio de transporte (según tecnología y modelo de operación seleccionados en la fase de diseño), reguladores (DGT, Comisión Europea), usuarios
(personal sanitario \textsc{chuf}), sindicatos sanitarios.
Matriz Poder--Interés con justificación por cuadrante.
\end{pendiente}